\section{A 239-state machine by delayed erasure}

We use the inherited state labels throughout this section.  The starting
state is 95, the instruction dispatcher is 96, and the separate halt state
is 254.  Renumbering the working states consecutively changes none of the
arguments below.  Let $M_{240}$ be the pinned rotation240 table, and let
$M_{239}$ be obtained by the following three row edits and one state
deletion:
\[
\begin{array}{c|c|c}
\text{row}&M_{240}&M_{239}\\\hline
147/0&0\mathrm{R}169&0\mathrm{R}80\\
80/0&0\mathrm{R}25&0\mathrm{L}106\\
91/1&0\mathrm{L}50&1\mathrm{L}167
\end{array}
\]
Delete both rows of state 169.  There is no other incoming edge to 169,
so these edits leave a complete binary table with 239 working states and
478 transition rows.  The source program, register permutation, encoding,
and initialization convention are unchanged.

\subsection{Why keeping a bit temporarily saves a state}

The essential change is the write in row $91/1$.  Instead of immediately
erasing the scanned 1, the new machine retains it and takes an existing
leftward route through states 167 and 151.  State 151 then distinguishes
two contexts.  An ordinary visit to 91 has the preceding word $011$;
a visit induced by the deleted selector state 169 has a 0 in the position
tested by 151.  The first context returns to the old erasure computation;
the second takes the existing row $173/1$ and reaches selector state 56.
All cells eventually have exactly their old values.  This is a change in
the timing of erasure, not an abstraction that discards tape contents.

\subsection{Structural context and the selector guard}

For a tape $T\colon\mathbb Z\to\{0,1\}$ and head position $h$, the old
machine satisfies the following context predicates:
\[
\begin{array}{c|l}
q&\text{predicate at state }q\\\hline
67&T(h-1)=0\\
80&T(h-2)=0,\quad T(h-1)=1\\
91&T(h-3)=0,\quad T(h-2)=T(h-1)=1\\
169&T(h-1)=0.
\end{array}
\]
These are ordinary inductive invariants, valid even for an arbitrary tape
started in state 95.  States 52 and 124 are the only predecessors of 67,
and both relevant rows write 0 and move right.  The only predecessor of
80 is $67/1$, the only predecessor of 91 is $80/1$, and the only predecessor
of 169 is $147/0$.  Inspection of those rows proves preservation; the
initial state 95 satisfies the predicates vacuously.

One additional fact uses the interpreter representation:
\begin{equation}\label{eq:selector169guard}
q=169\ \land\ T(h)=1\quad\Longrightarrow\quad T(h+1)=1.
\end{equation}
This is not asserted for arbitrary tapes.  It holds along the blank-start
computation, including every subsequently interpreted instruction.  The
separate guard certificate checks all such visits in the exact blank
initialization and in the complete unbounded symbolic instruction proof.
At each symbolic visit, the right-hand word starts with a run of 1s whose
affine length has constant term at least 2 and nonnegative coefficients.
Consequently the current cell and the next cell are both 1 for every
allowed register file, every additional header gap, and every finite
spare list.  The check introduces no new symbolic rewrite rule.  Unary
self-loop induction cannot hide a visit to state 169, and the four
finite-list schemas are rederived from the literal table.
More precisely, the check finds 32 symbolic visits, all at increment-4
sites, with right-hand run length $G+K$, where $G\geq0$ is the additional
header gap and $21\leq K\leq59$.  The exact initialization has six further
visits, each checked literally.  The required two-cell margin therefore
has a generous uniform reserve; it is not an extrapolation from small
register values.

The direct complete runtime proof for $M_{239}$ supplies a second route:
it checks all 8,386 original guarded instruction cases, without using the
source interval restrictions employed during discovery.  Thus the new
machine does not depend on specializing the source program or dropping
any of its guard cases.

\subsection{Exact positive-block comparison}

Let $\phi(169)=80$ and $\phi(q)=q$ for every other control state.  The
configuration map preserves the complete tape and head position.  Every
old transition outside states 80, 91, and 169 agrees with the corresponding
new transition under $\phi$.  The following six comparisons cover the
exceptional configurations:
\[
\begin{array}{c|l|c|c|c}
\text{old state/read}&\text{additional tape facts}&
\text{old steps}&\text{new steps}&\text{common endpoint}\\\hline
80/0&T(h-1)=1&1&5&(25,h+1)\\
80/1&\text{none}&1&1&(91,h+1)\\
169/0&T(h-1)=0&1&3&(57,h+1)\\
169/1&T(h-1)=0,\ T(h+1)=1&1&5&(56,h+1)\\
91/0&\text{none}&1&1&(106,h+1)\\
91/1&T[h-3,h)=011&3&11&(2,h-3).
\end{array}
\]
The notation for the endpoint suppresses the tape only to keep the table
compact: in every row the entire old and new tapes are equal, including
both infinite zero exteriors when those are present.  The nontrivial new
control paths are
\begin{align*}
80/0 &: 80\to106\to92\to107\to125\to25,\\
169/0 &:80\to106\to93\to57,\\
169/1 &:80\to91\to167\to151\to173\to56,\\
91/1 &:91\to167\to151\to24\to36\to54\to71
       \to82\to94\to50\to20\to2.
\end{align*}
In the last row the old path is $91\to50\to20\to2$.  Both paths erase
exactly the three 1s at $h,h-1,h-2$ and leave every other cell unchanged.
The selector row $169/1$ erases exactly the original cell $h$, as the old
row does.  All other listed comparisons preserve the initial tape.

The checker \texttt{verify\_local239.py} executes these paths with every
unmentioned tape cell represented by an independent symbolic variable.
It rejects any attempt to read such a variable without an explicit
assumption.  This establishes the comparisons for arbitrary exterior
tapes; it is not a finite sample of concrete tapes.

\subsection{Halting equivalence and the consistency statement}

All comparison blocks have positive length: at most 3 old steps and at
most 11 new steps.  Beginning with corresponding configurations, iterate
the appropriate comparison block along the deterministic old run.  Its
endpoint is again a reachable old configuration and the corresponding
new configuration, so the invariant and guard remain available.  If the
old computation is infinite, positivity gives an infinite new computation.
If the old computation halts, finitely many positive old blocks reach its
halt configuration, and the matching new blocks reach the same halt state.
No comparison path halts prematurely.  Therefore
\[
M_{239}\downarrow\quad\Longleftrightarrow\quad M_{240}\downarrow.
\]
The blank initial configurations correspond because the state map fixes
95 and does not alter the head or tape.  Combining this reduction with the
inherited rotation240 source calibration gives the requested equivalence
between nonhalting of $M_{239}$ and $\operatorname{Con}(\mathrm{ZFC})$.
This is an ordinary mathematical proof with executable, whole-domain
evidence.  A completed Lean or Coq theorem for the exact 239-state table
is not claimed by this packet.

The fastest of the three discovered row-$80/0$ variants is the one above.
Its blank initialization takes 409,835 steps, an increase of 144 over
rotation240.  The exact full tape, head, and dispatcher state at that
point agree with the inherited initialized frame.  The longer initialization
is immaterial to the state-count reduction or halting equivalence.
There is also an exact clock identity.  If $n_{q,b}$ counts visits to row
$q/b$ in the old run between corresponding endpoints, then
\[
t_{239}-t_{240}
=4n_{80,0}+2n_{169,0}+4n_{169,1}+8n_{91,1}.
\]
The old initialization has counts $18,12,6,3$ in that order, giving
$4\cdot18+2\cdot12+4\cdot6+8\cdot3=144$.  Each new comparison block has
between one and five times the duration of its old block, so any finite
halting run also satisfies $t_{240}\leq t_{239}\leq5t_{240}$.
